The matrix form of $\iota_k$ is Fact \claimref{transport.rule-1-morita-trivial} applied with base $T_k$: the new plane's ladder pair supplies the matrix units, and the base embeds as $a \mapsto \operatorname{diag}(a, \hat a)$, the grade involution taken in $T_k$. Note the parity is stable: $\dim$ of floor $k$ is $n + 2k$, of the same parity as $n$, so ``even charge'' is a property of the whole tower.

Even charge: $n$ even makes $\Cl(p,q)$ (and every floor) central simple over $\R$ (its centre is $\R$, \claimref{primer.centre}) \cite{Porteous1995, Lounesto2001}. The grade involution is an $\R$-algebra automorphism fixing the centre $\R$ pointwise, so its innerness is an instance of the Skolem-Noether theorem \cite{Pierce1982}; nothing need be imported, because the implementing element is exhibited directly: the volume element $u = x_1 \cdots x_n$ of an orthogonal generator frame satisfies $u x = (-1)^{n-1} x u = -xu$ for each generator $x$ (\claimref{std-algebra.pseudoscalar-commutation}) (move $x$ through $n$ letters, one coincidence), and $u^2 = \pm 1$ is invertible; so $\Ad_u(x) = u x u^{-1} = -x = \hat x$ on generators, and two automorphisms agreeing on generators agree. The displayed conjugation is then the block computation $U \operatorname{diag}(a,a)\, U^{-1} = \operatorname{diag}(a, u a u^{-1})$ with $U = \operatorname{diag}(1, u)$, invertible in $\Mat_2(T_k)$.

Odd charge: the centre of each floor is $\R \oplus \R$ ($d \equiv 1 \bmod 4$) or $\C$ ($d \equiv 3 \bmod 4$), generated by $1$ and the volume element $\omega$, which is odd (its length $n + 2k$ is odd) and is therefore sent by the grade involution to $-\omega$: the involution exchanges the two central idempotents $\tfrac12(1 \pm \omega)$, respectively conjugates the central complex structure. An inner automorphism fixes the centre pointwise; the grade involution does not, so it is outer, and no unitary correction of the diagonal exists.
